\section{Estimation method}
\label{sec:estimation}

The estimator uses the observed data and the centering conventions of Section~\ref{sec:model}. Write $q=d+K$, $P=p+q$, $\bm{J}=\bm{I}_n-\mathbf 1_n\mathbf 1_n^\top/n$, $\bm{X}_c=\bm{J}\bm{X}$, and $\bm{Z}_c=\bm{J}\bm{Z}$. Its primary target is the sparse molecular slope in $\mathcal M_0$, with $\mathcal M_1$ interpreted relative to its specified latent representation.

\subsection{Covariate-only factor estimation}
\label{subsec:pca-estimator}
Take $n\ge2$ and a prespecified $K\le\operatorname{rank}(\bm{X}_c)\le n-1$. For $K>0$, let $\bm{v}_1,\ldots,\bm{v}_K$ be orthonormal eigenvectors of $\bm{X}_c\bm{X}_c^\top/n$ associated with its largest $K$ eigenvalues $\widehat\ell_1\ge\cdots\ge\widehat\ell_K>0$, and define
\begin{equation}
 \begin{aligned}
 \widehat{\bm{F}}&=\sqrt n\,(\bm{v}_1,\ldots,\bm{v}_K),&
 \widehat{\bm{\Lambda}}&=\operatorname{diag}(\widehat\ell_1,\ldots,\widehat\ell_K),\\
 \widehat{\bm{B}}&=\bm{X}_c^\top\widehat{\bm{F}}/n,&
 \widehat{\bm{U}}&=\bm{X}_c-\widehat{\bm{F}}\widehat{\bm{B}}^\top.
 \end{aligned}
 \label{eq:pca-construction}
\end{equation}
Signs and tied eigenspaces follow a deterministic covariate-measurable convention. A tie across the $K$th eigenvalue can alter the selected subspace; the recovery analysis assumes spectral separation at $K$. The construction is outcome-free and gives
\begin{equation}
 \begin{gathered}
 \widehat{\bm{F}}^\top\widehat{\bm{F}}/n=\bm{I}_K,\qquad
 \widehat{\bm{F}}^\top\widehat{\bm{U}}=0,\qquad
 \bm{J}\widehat{\bm{F}}=\widehat{\bm{F}},\quad \bm{J}\widehat{\bm{U}}=\widehat{\bm{U}},\\
 \widehat{\bm{B}}^\top\widehat{\bm{B}}=\widehat{\bm{\Lambda}},\qquad
 \bm{X}_c\widehat{\bm{B}}=\widehat{\bm{F}}\widehat{\bm{\Lambda}}
 \end{gathered}
 \label{eq:pca-identities}
\end{equation}
directly. This principal-component residualization follows \citet{FanEtAl2020}; sample factor--residual orthogonality allows dependence among the residual coordinates.

Let $\widehat{\bm{f}}_i$ and $\widehat{\bm{u}}_i$ denote the column vectors corresponding to row $i$ of $\widehat{\bm{F}}$ and $\widehat{\bm{U}}$. Set
\begin{equation}
 \begin{gathered}
\widehat{\bm{A}}=(\bm{Z}_c,\widehat{\bm{F}}),\\
 \widehat{\bm{W}}=(\widehat{\bm{A}},\widehat{\bm{U}}),\\
 \bm{a}=(\bm{\theta}^\top,\bm{\gamma}^\top)^\top,\quad
 \bm{\vartheta}=(\bm{a}^\top,\bm{\beta}^\top)^\top.
 \end{gathered}
 \label{eq:feasible-design}
\end{equation}
Thus $\widehat{\bm{W}}$ is $n\times P$, with row $\widehat{\bm{W}}_i^\top$. The convention $K=0$ omits the factor block and sets $\widehat{\bm{U}}=\bm{X}_c$, defining the unadjusted comparison method.

Residual coefficients retain the molecular coordinate scales fixed in Section~\ref{sec:model}. Under a rescaling $\widetilde{\bm{U}}=\widehat{\bm{U}} \bm{S}^{-1}$ with positive diagonal $\bm{S}=\operatorname{diag}(s_1,\ldots,s_p)$, the equivalent coordinates and penalty are $\bm{\zeta}=\bm{S}\bm{\beta}$ and $\lambda\sum_j|\zeta_j|/s_j$. All learned transformations are confined to the training sample.

\subsection{Censored pairwise loss and scalar smoothing}
\label{subsec:smoothed-loss}
For a finite parameter vector, define
\begin{equation}
 \begin{gathered}
e_i(\bm{\vartheta})=Y_i-\widehat{\bm{W}}_i^\top\bm{\vartheta},\\
 r_{ij}(\bm{\vartheta})=e_j(\bm{\vartheta})-e_i(\bm{\vartheta}),\\
 \bm{w}_{ij}=\widehat{\bm{W}}_i-\widehat{\bm{W}}_j,\\
 N=n(n-1).
 \end{gathered}
 \label{eq:residual-pairs}
\end{equation}
The unsmoothed Gehan loss is
\begin{equation}
 L_{n,0}(\bm{\vartheta})
 =\frac1N\sum_{i\ne j}\Delta_i\{r_{ij}(\bm{\vartheta})\}_+.
 \label{eq:gehan-feasible}
\end{equation}
The criterion follows regularized Gehan estimation \citep{CaiEtAl2009,SuderMolstad2022}. Event patients anchor the index $i$, and every other patient, including censored patients, contributes as $j$. We retain the ordered off-diagonal normalization $N=n(n-1)$ throughout.

Let $V\sim N(0,1)$ be an auxiliary integration variable, unrelated to the AFT error. For a common scalar bandwidth $h>0$, define
\begin{equation}
 \begin{gathered}
\ell_h(r)=E_V(r+hV)_+
   =r\Phi(r/h)+h\phi(r/h),\\
 L_{n,h}(\bm{\vartheta})=\frac1N\sum_{i\ne j}\Delta_i\ell_h\{r_{ij}(\bm{\vartheta})\},
 \end{gathered}
 \label{eq:smoothed-gehan}
\end{equation}
where $\phi$ and $\Phi$ are the standard normal density and distribution function. The Gaussian convolution is evaluated analytically while the AFT error distribution remains unspecified. Differentiation gives
\begin{equation}
 \ell_h'(r)=\Phi(r/h),\qquad
 \ell_h''(r)=h^{-1}\phi(r/h)>0.
 \label{eq:scalar-derivatives}
\end{equation}
Consequently $\ell_h$ is convex and 1-Lipschitz, with derivative $1/2$ at zero.

\begin{proposition}[Uniform loss approximation]
\label{prop:smoothing-approximation}
For every $r\in\mathbb R$ and $h>0$,
\begin{equation}
 \begin{gathered}
0\le \ell_h(r)-r_+\le\frac{h}{\sqrt{2\pi}},\\
 \int_{\mathbb R}|\Phi(r/h)-\mathbf 1\{r\ge0\}|\,\mathrm{d}r
   =h\sqrt{\frac2\pi}.
 \end{gathered}
 \label{eq:smoothing-bounds}
\end{equation}
For any fixed observed sample, writing $\bar\Delta=n^{-1}\sum_i\Delta_i$,
\begin{equation}
 0\le L_{n,h}(\bm{\vartheta})-L_{n,0}(\bm{\vartheta})
   \le \frac{\bar\Delta h}{\sqrt{2\pi}}
 \quad\text{uniformly over }\bm{\vartheta}\in\mathbb R^P.
 \label{eq:uniform-loss-bound}
\end{equation}
\end{proposition}
Appendix~\ref{app:method-properties} proves the deterministic identities in this section. Proposition~\ref{prop:smoothing-approximation} controls the loss approximation; Section~\ref{sec:theory} combines bandwidth-dependent score bias, identification, and curvature to obtain coefficient bounds.

Scalar smoothing of censored rank scores has precedents in \citet{Heller2007}. It uses a common bandwidth, whereas the parameter-perturbation construction of \citet{JohnsonStrawderman2009} induces pair-specific bandwidths. For a perturbation $n^{-1/2}\bm{\Gamma} \bm{V}_P$, with $\bm{V}_P\sim N(0,\bm{I}_P)$, these are
\[
 h_{ij}=n^{-1/2}\{\bm{w}_{ij}^\top\bm{\Gamma}\bm{\Gamma}^\top \bm{w}_{ij}\}^{1/2},
\]
which generally varies by pair. The common scalar $h$ used here leads to the smoothing-bias analysis in Section~\ref{sec:theory}.

\subsection{Partial regularization and finite minimizers}
\label{subsec:partial-penalty}
For fixed $h>0$ and $\lambda>0$, let
\begin{equation}
 \begin{gathered}
\mathcal S_{n,h,\lambda}
 =\operatorname*{arg\,min}_{\vartheta\in\mathbb R^P}
 Q_{n,h,\lambda}(\bm{\vartheta}),\\
 Q_{n,h,\lambda}(\bm{\vartheta})=L_{n,h}(\bm{\vartheta})+\lambda\|\bm{\beta}\|_1.
 \end{gathered}
 \label{eq:estimator-set}
\end{equation}
When this set is nonempty, $\widehat{\bm{\vartheta}}$ denotes a member. The $\ell_1$ penalty produces sparsity in the $p$ molecular coefficients, and clinical and factor effects are estimated jointly with them. Section~\ref{sec:theory} establishes coefficient-error guarantees.

The identity $\widehat{\bm{U}}=\bm{X}_c-\widehat{\bm{F}}\widehat{\bm{B}}^\top$ implies
\begin{equation}
 \begin{gathered}
\widehat{\bm{W}}\bm{\vartheta}
 =\bm{Z}_c\bm{\theta}+\bm{X}_c\bm{\beta}+\widehat{\bm{F}}\bm{\nu},
 \\
 \bm{\nu}=\bm{\gamma}-\widehat{\bm{B}}^\top\bm{\beta}.
 \end{gathered}
 \label{eq:relaxed-predictor}
\end{equation}
For fixed training data,~\eqref{eq:relaxed-predictor} is a bijective reparameterization preserving $\|\bm{\beta}\|_1$. The observed-$\bm{X}$ formulation includes a freely estimated principal-component term. Setting $\bm{\gamma}=\widehat{\bm{B}}^\top\bm{\beta}$ recovers the ordinary penalized smoothed AFT fit; allowing $\bm{\gamma}$ to vary gives the proposed relaxation. Both blocks are optimized jointly under the nonlinear pairwise loss.

Finite attainment with the unpenalized block is governed by the nuisance-design geometry. A sufficient sample condition, when $q>0$, is
\begin{equation}
 \chi_n:=\inf_{\|\bm{v}\|_2=1}
 \frac1N\sum_{i\ne j}\Delta_i
       \{(\widehat{\bm{A}}_i-\widehat{\bm{A}}_j)^\top \bm{v}\}_+>0.
 \label{eq:sample-no-separation}
\end{equation}
Condition~\eqref{eq:sample-no-separation} makes $\mathcal S_{n,h,\lambda}$ nonempty and compact; Appendix~\ref{subapp:existence} proves this result and illustrates nonattainment. For $q=0$, the positive $\ell_1$ penalty bounds the sublevel sets. Section~\ref{sec:theory} also establishes finite attainment on the statistical event used for estimation bounds.

A deterministic tie-breaking rule may select among multiple minimizers; numerical eligibility is recorded separately.

\subsection{Reference coordinates and the approximation term}
\label{subsec:reference-coordinates}
To relate~\eqref{eq:estimator-set} to the target, write the centered latent matrices as $\bm{F}_c=\bm{J}\bm{F}$ and $\bm{U}_c=\bm{J}\bm{U}$. For any invertible alignment matrix $\bm{H}$, set
\begin{equation}
 \begin{gathered}
\bm{E}_F=\widehat{\bm{F}}-\bm{F}_c\bm{H}^\top,\\
 \bm{\gamma}^\dagger=\widehat{\bm{B}}^\top\bm{\beta}_0+\bm{H}^{-T}\bm{\xi}_0,\\
 \bm{\vartheta}^\dagger=(\bm{\theta}_0^\top,(\bm{\gamma}^\dagger)^\top,\bm{\beta}_0^\top)^\top.
 \end{gathered}
 \label{eq:reference-parameter}
\end{equation}
These reference coordinates link the feasible criterion to the population target. In $\mathcal M_0$, $\bm{\gamma}^\dagger=\widehat{\bm{B}}^\top\bm{\beta}_0$ and the alignment matrix drops out. Recall
\[
 \bm{\eta}^0=\bm{Z}\bm{\theta}_0+\bm{X}\bm{\beta}_0+\bm{F}\bm{\xi}_0,\qquad
 \bm{\eta}_c^0=\bm{J}\bm{\eta}^0,\qquad
 R_i=Y_i-\eta_i^0.
\]

\begin{proposition}[Exact reference representation]
\label{prop:reference-representation}
With $\mathbf d=\bm{E}_F\bm{H}^{-T}\bm{\xi}_0$ and $b_n=\|\mathbf d\|_\infty$,
\begin{equation}
 \begin{aligned}
 \widehat{\bm{W}}\bm{\vartheta}^\dagger&=\bm{\eta}_c^0+\mathbf d, & \bm{J}\mathbf d&=\mathbf d,\\
 e_i(\bm{\vartheta}^\dagger)&=R_i+\bar\eta^0-d_i, &
 r_{ij}(\bm{\vartheta}^\dagger)&=R_j-R_i+d_i-d_j.
 \end{aligned}
 \label{eq:reference-identities}
\end{equation}
Moreover,
\begin{equation}
 \begin{gathered}
b_n\le
 \left(\max_i\|(\bm{E}_F)_{i\cdot}\|_2\right)\|\bm{H}^{-T}\bm{\xi}_0\|_2,
 \\
 \widehat{\bm{W}}\widehat{\bm{\vartheta}}-\bm{\eta}_c^0
 =\widehat{\bm{W}}(\widehat{\bm{\vartheta}}-\bm{\vartheta}^\dagger)+\mathbf d.
 \end{gathered}
 \label{eq:approximation-decomposition}
\end{equation}
In $\mathcal M_0$, $\mathbf d=0$ exactly for every admissible PCA output.
\end{proposition}
The common location in~\eqref{eq:reference-identities} cancels within each pair. Thus $b_n$ measures approximation of the additional factor contribution. For $\mathcal M_0$, the ordinary predictor is represented exactly and factor estimation enters through the estimated-design geometry. Nuisance error is measured from the sample-dependent reference $\bm{\gamma}^\dagger$, allowing $\bm{B}^\top\bm{\beta}_0$ to grow with the sparse signal.

Within a fixed estimated subspace, replacing $\widehat{\bm{F}}$ and $\widehat{\bm{B}}$ by $\widehat{\bm{F}}\bm{Q}^\top$ and $\widehat{\bm{B}}\bm{Q}^\top$, and $\bm{\gamma}$ by $\bm{Q}\bm{\gamma}$, preserves $\widehat{\bm{U}}$, the predictor, and the objective. The fitted-predictor and molecular-minimizer sets are invariant under this coordinate rotation. Subspace selection at a spectral tie is governed by the convention in Section~\ref{subsec:pca-estimator}.

\subsection{Training-frozen scores for held-out patients}
\label{subsec:out-of-sample-map}
For a new patient, the training means and loadings define a frozen prediction map. Let $\bm{x}_c=\bm{x}-\bar{\bm{X}}$ and $\bm{z}_c=\bm{z}-\bar{\bm{Z}}$. For $K>0$, set
\begin{equation}
 \begin{gathered}
\widehat{\bm{f}}(\bm{x})=\widehat{\bm{\Lambda}}^{-1}\widehat{\bm{B}}^\top \bm{x}_c,\\
 \widehat{\bm{u}}(\bm{x})=\bm{x}_c-\widehat{\bm{B}}\widehat{\bm{f}}(\bm{x}),\\
 \widehat\eta(\bm{x},\bm{z})=\bm{z}_c^\top\widehat{\bm{\theta}}+
       \widehat{\bm{f}}(\bm{x})^\top\widehat{\bm{\gamma}}+
       \widehat{\bm{u}}(\bm{x})^\top\widehat{\bm{\beta}}.
 \end{gathered}
 \label{eq:frozen-score-map}
\end{equation}
For $K=0$, omit the factor terms. By~\eqref{eq:pca-identities}, this least-squares projection onto the estimated loading space reproduces the fitted training factors and residuals exactly.

Equivalently,
\begin{equation}
 \begin{gathered}
\widehat\eta(\bm{x},\bm{z})=\bm{z}_c^\top\widehat{\bm{\theta}}+\bm{x}_c^\top\widehat{\bm{b}}_{\rm pred},\\
 \widehat{\bm{b}}_{\rm pred}
 =\widehat{\bm{\beta}}+\widehat{\bm{B}}\widehat{\bm{\Lambda}}^{-1}
         (\widehat{\bm{\gamma}}-\widehat{\bm{B}}^\top\widehat{\bm{\beta}}).
 \end{gathered}
 \label{eq:prediction-coefficients}
\end{equation}
The sparse coefficient $\widehat{\bm{\beta}}$ summarizes residual effects, whereas $\widehat{\bm{b}}_{\rm pred}$ combines residual and factor-mediated contributions and can be dense. Higher $\widehat\eta$ denotes longer conditional event times; $-\widehat\eta$ supplies the opposite risk ordering. The score represents relative log-time location. Calibrated event-time quantiles additionally require the error distribution and its location.

\subsection{Bandwidth, penalty, and training selection}
\label{subsec:tuning-interface}
The estimator is indexed by $(K,h,\lambda)$. The coefficient analysis takes the correct $K$ as fixed, with $h=h_n>0$ and $\lambda=\lambda_n>0$ deterministic or fixed after conditioning on the design. Section~\ref{sec:theory} relates the smoothing bias, $b_n$, and stochastic fluctuations to the penalty and gives sufficient bandwidth and penalty conditions for the stated error rates.

For empirical tuning, let $\mathcal V_1,\ldots,\mathcal V_L$ be held-out folds of sizes $m_\ell\ge2$. For every candidate $(K,h,\lambda)$, all learned preprocessing, PCA, and coefficient fitting use the complement of $\mathcal V_\ell$. Apply~\eqref{eq:frozen-score-map} to obtain $\widehat\eta_i^{(-\ell)}$ and put $e_i^{(-\ell)}=Y_i-\widehat\eta_i^{(-\ell)}$. The cross-validation (CV) score is
\begin{equation}
 \operatorname{CV}(K,h,\lambda)=
 \frac{\displaystyle\sum_{\ell=1}^L
       \sum_{\substack{i,j\in\mathcal V_\ell\\ i\ne j}}
       \Delta_i\left[\{e_j^{(-\ell)}-e_i^{(-\ell)}\}_+
                       -\{Y_j-Y_i\}_+\right]}
      {\displaystyle\sum_{\ell=1}^Lm_\ell(m_\ell-1)}.
 \label{eq:within-fold-cv}
\end{equation}
The candidate-independent baseline subtraction preserves the selection while expressing validation as a finite loss contrast. The same unsmoothed score is used for every candidate bandwidth.

Within-fold differences remove arbitrary fold-specific locations. All candidate bandwidths share the same unsmoothed score, computed over within-fold pairs. Alternative pooling rules are studied by \citet{SuderMolstad2022}. Section~\ref{subsec:cv-results} specifies numerical eligibility and tie handling for the empirical study. The coefficient theorems use fixed or covariate-measurable tuning; the outcome-selected procedure is evaluated separately.

\subsection{Derivative and optimality interface}
\label{subsec:optimality-interface}
For $h>0$, differentiation of the finite sum gives
\begin{equation}
 \begin{aligned}
 \nabla L_{n,h}(\bm{\vartheta})
 &=\frac1N\sum_{i\ne j}\Delta_i\Phi\{r_{ij}(\bm{\vartheta})/h\}\bm{w}_{ij},\\
 \nabla^2L_{n,h}(\bm{\vartheta})
 &=\frac1{Nh}\sum_{i\ne j}\Delta_i\phi\{r_{ij}(\bm{\vartheta})/h\}\bm{w}_{ij}\bm{w}_{ij}^\top.
 \end{aligned}
 \label{eq:loss-derivatives}
\end{equation}
The positive-semidefinite Hessian establishes convexity. A finite vector is a global minimizer of~\eqref{eq:estimator-set} if and only if
\begin{equation}
 \begin{gathered}
\nabla_aL_{n,h}(\widehat{\bm{\vartheta}})=0,\\
 0\in\nabla_\beta L_{n,h}(\widehat{\bm{\vartheta}})
          +\lambda\partial\|\widehat{\bm{\beta}}\|_1.
 \end{gathered}
 \label{eq:kkt-interface}
\end{equation}
A nonzero molecular coordinate has gradient $-\lambda\operatorname{sign}(\widehat\beta_j)$, and a zero coordinate has gradient in $[-\lambda,\lambda]$. Numerical stationarity is assessed on these coordinates and the complete unpenalized block.

The gradient admits an efficient matrix form. Define $V_{ii}=0$, $V_{ij}=\Delta_i\Phi(r_{ij}/h)$ for $i\ne j$, and $c_i=\sum_jV_{ij}-\sum_jV_{ji}$. Then
\begin{equation}
 \nabla L_{n,h}(\bm{\vartheta})=\widehat{\bm{W}}^\top \bm{c}/N.
 \label{eq:matrix-gradient}
\end{equation}
